\subsection{Líneas de Investigación Futura}
\label{subsec:lineas_investigacion_futura}

A partir de los hallazgos y limitaciones identificadas, se proponen las siguientes líneas de trabajo para la comunidad investigadora:

\begin{enumerate}
    \item \textbf{Canalizaciones Completamente Locales y Abiertas (\textit{Open-Weights}):} Sustituir las dependencias de APIs en la nube por modelos fundacionales multimodales abiertos ejecutables en local, tales como la familia \textit{Qwen3-VL-Embedding} y \textit{Qwen3-VL-Reranker} \cite{qwen3_vl_emb2026}, o arquitecturas ultraligeras de visión-lenguaje como \textit{PaliGemma} y \textit{Moondream}, garantizando privacidad total y soberanía de datos en entornos empresariales.
    \item \textbf{Video-RAG en Tiempo Real sobre Flujos Continuos (\textit{Streaming}):} Adaptar los algoritmos de segmentación SSIM y alineación Whisper a ventanas deslizantes de memoria en tiempo real, permitiendo indexar y consultar transmisiones de vídeo en vivo (conferencias, emisiones deportivas o vigilancia de seguridad) con latencias de respuesta inferiores a un segundo.
    \item \textbf{Ponderación Modal Aprendida (\textit{Attention-Guided Fusion}):} Sustituir la concatenación estática de partes visuales y textuales por un mecanismo de atención cruzada que pondere dinámicamente el peso relativo de la imagen frente a la voz en función de la correlación semántica detectada en cada plano.
    \item \textbf{Cuantización Extrema sobre Representaciones Matryoshka:} Combinar el truncamiento MRL con técnicas de cuantización escalar a 8 bits (\texttt{int8}) o binarización vectorial a 1 bit, proyectando reducciones adicionales de memoria de hasta $32\times$, lo que posibilitaría mantener miles de horas de vídeo indexadas en dispositivos móviles o de computación en el borde (\textit{edge computing}).
\end{enumerate}
